\chapterstart{Обучение и правила сравнительного исследования}
\label{sec:design}
\subsection{Обновление моделей и классическая обратная связь}
Обучение чередует сбор переходов среды и обновление нейросетей по буферу опыта. Одни и те же переходы могут использоваться повторно, поэтому методы называют \emph{off-policy}: текущая политика учится по опыту, собранному в том числе прежними версиями. Критик приближает дисконтированную ценность запроса, а актор, если он есть, формирует запрос по наблюдению. Целевая сеть — медленно меняющаяся копия, используемая для вычисления ориентиров обучения.

Ниже $o,a,r,o'$ обозначают наблюдение, запрос, награду и следующее наблюдение; $d=1$ означает терминальное событие, без продолжения оценки будущего возврата. Временное усечение по горизонту этим событием не считается. $Q_\psi$ — критик с параметрами $\psi$, $\mu_\eta$ или $\pi_\eta$ — актор с параметрами $\eta$; черта над параметрами обозначает целевую сеть. Величина $y$ — целевое значение для обновления критика, а не непосредственно измеренный будущий возврат.

Для DQN цель имеет вид $y=r+\gamma(1-d)\max_{a'}Q_{\bar\psi}(o',a')$. Сеть приближает эту цель по сохранённым переходам; при оценке выбирается индекс максимальной ценности. В обучении применяется $\varepsilon$-жадное исследование: вероятность случайного запроса линейно уменьшается от 1 до 0,05 за первые 20\% бюджета. Целевая сеть обновляется каждые 1000 переходов. Рассматривается обычный DQN из используемой библиотеки; дополнительные алгоритмы под тем же именем не вводятся.

DDPG использует цель $y=r+\gamma(1-d)Q_{\bar\psi}(o',\mu_{\bar\eta}(o'))$, где $\mu_\eta$ — детерминированный актор. Его обновление направлено на увеличение $Q_\psi(o,\mu_\eta(o))$. В конфигурации один критик; исследование обеспечивается независимым гауссовым шумом со стандартным отклонением $\sigma=0{,}1$ в нормированных запросах. Шум отключён при оценке. Реализация использует ровно один критик: двойной критик и иные модификации базового DDPG здесь не включены.

SAC оптимизирует не только возврат, но и энтропию запросов. Для выбранного следующего запроса $a'\sim\pi_\eta(\cdot|o')$ цель двух критиков записывается как
\begin{equation}
y=r+\gamma(1-d)\left[\min_{j=1,2}Q_{\bar\psi_j}(o',a')-\alpha\log\pi_\eta(a'|o')\right].
\label{eq:sac}
\end{equation}
Актор минимизирует среднее $\alpha\log\pi_\eta(a|o)-\min_jQ_{\psi_j}(o,a)$. Коэффициент $\alpha$ в~\eqref{eq:sac} настраивается автоматически с начальным значением 0,1. Внешняя оценка использует детерминированный режим библиотечной политики, а не новый случайный выбор действий. Это важно: случайность обучения и случайность исполнения не смешиваются.

TQC сохраняет энтропийного актора, но вместо одного числа каждый критик предсказывает 25 квантилей распределения возврата. В исследованной конфигурации два критика, параметр удаления верхних квантилей равен двум на сеть: из объединённых 50 прогнозов для цели остаются 46 нижних. Обучение критиков использует квантильную регрессию. Этот механизм предназначен для контроля переоценивания, однако эксперимент здесь проверяет итоговую процедуру TQC целиком, без отдельного выключения усечения.

Классический контроллер сначала следует сохранённому решению прямой коллокации — поиску состояний и входов на временной сетке с ограничениями динамики: номинальным состоянию $q_0(t)$ и ускорению $u_0(t)$. TVLQR (Time-Varying Linear Quadratic Regulator) — линейно-квадратичный регулятор с зависящими от времени коэффициентами — задаёт запрос
\begin{equation}
u_{\rm req}(t)=u_0(t)-K(t)(q-q_0(t))-k_0(t),\qquad q=(x,\theta,v,\omega).
\label{eq:tracking}
\end{equation}
Аффинный член $k_0$ в~\eqref{eq:tracking} учитывает дефект восстановленной номинальной траектории. После конца плана выполняется однократный переход к стабилизирующему LQR: $u=-K_b(x,e_\theta,v,\omega)^\mathrm{T}$, с матрицей штрафа за состояние $Q=\operatorname{diag}(10,1,10,1)$ и штрафом за управление $R=0{,}5$. Ошибка угла периодическая только в режиме балансировки; слежение сохраняет запланированный полный оборот. Повторного планирования и нового подъёма после неудачного захвата нет. Таким образом, это конкретный классический baseline, а не оптимальный регулятор для каждого нового старта.

\subsection{Общие условия и единицы сравнения}
Исходное распределение сброса среды (\emph{reset}) — произведение независимых равномерных распределений:
\begin{equation}
\begin{aligned}
x_0&\sim U[-0{,}02,0{,}02]\ \text{м}, &v_0&\sim U[-0{,}02,0{,}02]\ \text{м/с},\\
\theta_0&\sim U[-0{,}05,0{,}05]\ \text{рад}, &\omega_0&\sim U[-0{,}05,0{,}05]\ \text{рад/с}.
\end{aligned}\label{eq:reset}
\end{equation}
Значение seed обучения задаёт поток случайности, но не границы~\eqref{eq:reset}. Для on и off весь этот малый прямоугольник допустим при сбросе. Посещение других состояний во время движения не означает, что они использовались как исходные состояния обучения.

Каждое обучение в анализируемых сериях содержит 100\,000 переходов, одну среду, буфер на 100\,000 записей и пакеты по 64 перехода (размер пакета, batch size). Сети имеют два скрытых слоя по 64 нейрона с ReLU, функцией $z\mapsto\max(0,z)$; параметры обновляет оптимизатор Adam по оценке градиента из пакета опыта. Базовая скорость обучения равна $3\cdot10^{-4}$. После 128 начальных переходов выполняется одно обновление на каждый переход, всего 99\,872. У непрерывных методов коэффициент мягкого обновления целей равен 0,005; у DQN действуют указанные выше периодические обновления. Физика, интегратор, R1 и семантика запросов общие (Приложение~\ref{app:repro}).

Фиксированный бюджет сравнивает процедуры при одинаковом числе взаимодействий, но не при одинаковом вычислительном времени и не при достигнутой сходимости. У TQC дополнительные квантильные прогнозы меняют объём вычислений; равенство числа переходов не задаёт равенства стоимости обучения. В сериях использовалось параллельное исполнение независимых обучений; сумма длительностей процессов не является временем всей серии. Системные диагностические пилоты и замеры производительности не входят в 33 обучения Таблицы~\ref{tab:series}.

\begin{table}[htbp]\caption{Состав анализируемых серий и назначение данных}\label{tab:series}
\begin{smalltable}\begin{tabularx}{\linewidth}{@{}p{22mm} X c p{37mm}@{}}\toprule
Серия & Что сравнивается & Обучений & Оценочные старты\\\midrule
Основная & SAC/TQC/DDPG/DQN с фильтром; сопоставимый SAC без него, seed 0–2 & 15 & 20 валидационных + 3 именованных\\
Stage 1 & Два изменения DDPG-on и SAC-off, seeds 0–2 & 12 & Прежние 20 + 3\\
Confirmation & Фиксированные warmup5000, новые seeds 3–5 & 6 & Прежние 20 + 3\\
Robustness & Шесть лучших исходных политик и классика, исполнение off/on & 0 & 140 новых: 7 групп\\\bottomrule
\end{tabularx}\end{smalltable}\end{table}

\subsection{Выбор контрольной точки и проверка настройки}
Через каждые 10\,000 обучающих переходов, включая исходную и конечную модели, проводится детерминированная оценка на 20 фиксированных validation-состояниях из~\eqref{eq:reset}. Ещё три именованных сценария (\emph{named}) задают точное нижнее положение и $\theta_0=\pm0{,}02$ при остальных нулевых координатах. Они описательные и никогда не входят в выбор.

Лучшей точкой (\emph{best}) называется сохранённая модель, максимизирующая лексикографический ключ: сначала сравнивается первая компонента, при равенстве — следующая
\begin{equation}
K_{\rm ckpt}=(p,\overline H_f,\overline H_\ell,-\overline t_{\rm first},n).
\label{eq:best}
\end{equation}
В~\eqref{eq:best} $p$ — доля успеха на validation20, черта означает среднее, $t_{\rm first}$ — начало первого подтверждённого интервала удержания длиной не менее 2 с (среднее по эпизодам, где он есть), $n$ — число обучающих переходов. Если такого интервала нет ни в одном эпизоде, вместо среднего времени используется 11 с. При равных физических метриках выбирается более поздняя контрольная точка. Возврат в ключ не входит. Последней точкой (\emph{last}) называются веса после 100\,000 переходов. Best и last иногда совпадают; их два названия не создают два независимых обучения.

Внешняя оценка сохранённых best/last выполняется при обоих режимах фильтра на тех же 23 стартах: $2\times2\times23=92$ слота на обучение. Классика имеет 46. Это проверка сохранённых моделей и вмешательства в их исполнение, но не независимая оценка обобщения: validation уже использована в выборе.

Stage 1 меняет ровно одно значение относительно baseline: либо начальный сбор опыта 128 $\to$ 5000 (\texttt{warmup5000}), либо скорость обучения $3\cdot10^{-4}\to10^{-4}$ (\texttt{lr1e4}). Во втором случае остаётся 99\,872 обновления, в первом — 95\,000. Смена warmup затрагивает последовательность опыта и число обновлений, поэтому одна изменённая настройка не означает один изолированный причинный механизм.

Настройка выбирается по трём seeds и только валидации лучшей точки в режиме соответствующего обучения:
\begin{equation}
K_{\rm cfg}=(\min_i p_i,\operatorname{mean}_i p_i,
\min_i\overline H_{f,i},\operatorname{mean}_i\overline H_{f,i},
\operatorname{mean}_i\overline H_{\ell,i}).
\label{eq:config}
\end{equation}
По~\eqref{eq:config} требуется строгое улучшение исходной настройки; при равенстве сохраняется базовый вариант. Для SAC-off используется оценка off, а улучшение от включения фильтра при оценке в этот ключ не подставляется.

Выбранные настройки затем без изменений проверяются на новых значениях seed 3, 4, 5. Предварительно закреплённый критерий confirmation: каждый новый best имеет $p\geq0{,}8$ (не менее 16/20), не менее двух из трёх last также имеют $p\geq0{,}8$, все задания и оценки завершены. Это правило принятия проекта, а не статистический тест с заданной вероятностью ошибки. Новые значения seed проверяют повторяемость обучения; старты оценки остаются прежними и не становятся независимыми от отбора.

\subsection{Что проверяют данные и что разработано в проекте}
Для обучения сохранены конфигурации, веса, replay/runtime, счётчики и происхождение checkpoint; для оценки — полные журналы состояний и переходов, идентичности стартов и контрольные суммы. Завершённые аудиты отдельно проверили основную серию, Stage 1 и confirmation (Приложение~\ref{app:repro}). Подготовка настоящего отчёта использует их компактные таблицы и адресно сверенные сохранённые примеры; полный аудит не повторяется.

Уровни свидетельств различаются. SHA выявляет изменение байтов. Согласованность компактных таблиц проверяет идентичности и арифметику. Пересчёт полных журналов проверяет определения метрик по сохранённым переходам, но не повторно интегрирует систему. Ни один из этих уровней не является испытанием оборудования. Полные высокоточные обучающие траектории сохранены не для всех переходов. Имеются накопленные статистики и отдельные контрольные примеры; они не позволяют восстановить отсутствующий временной ряд без предположений.

По истории репозитория в проекте добавлены исправленный интерфейс симулятора, слой событий рабочих ограничений, Gymnasium-обёртка, общий фильтр и его подключение, журналы, повторяемое исполнение и анализ. Лабораторная динамика, базовые LQR/траекторные идеи и алгоритмы библиотек имеют самостоятельное происхождение. Эти результаты представлены как выполненная в рамках курсовой инженерная и аналитическая работа с существенной помощью ИИ, а не как единолично написанные с нуля физический симулятор или алгоритмы обучения. Подробное раскрытие приведено в Приложении~\ref{app:repro}.

Итог протокола — сравнение разных источников изменчивости без смешения знаменателей. Эпизод, уникальный физический старт, checkpoint и независимый запуск обучения отвечают на разные вопросы. Следующие главы используют эту структуру при интерпретации результатов.
